\documentclass[10pt,a4paper]{amsart}
\usepackage[margin=19mm]{geometry}
\usepackage{amsmath,amssymb,mathtools}
\usepackage[scr=rsfs,cal=euler]{mathalfa}
\usepackage{xfrac}
\usepackage[unicode,hidelinks,bookmarks=false]{hyperref}
\newcommand{\ie}{i.e.\ }
\newcommand{\cf}{cf.\ }
\newcommand{\resp}{respectively\ }
\newcommand{\Subsection}{\subsection}
\providecommand{\nobreakdash}{\nobreak\hbox{-}}
\setlength{\emergencystretch}{2em}
\multlinegap0pt
\usepackage{amssymb}
\newtheorem{theorem}{Theorem}
\newcommand{\C}{{\mathbb{C}}}
\renewcommand{\P}{{\mathbb{P}}}
\title{Total curvature of a complete minimal surface and the modified defect relation of a Fermat hypersurface for the Gauss map}
\author{Si Duc Quang, Nguyen Thi Quynh Chi}
\date{Source: arXiv:2605.27001v1 (CC BY 4.0)}
\begin{document}
\maketitle

\noindent\emph{Source and licence.} Total curvature of a complete minimal surface and the modified defect relation of a Fermat hypersurface for the Gauss map. Version arXiv:2605.27001v1 (CC BY 4.0), distributed by arXiv under the Creative Commons Attribution 4.0 International licence, http://creativecommons.org/licenses/by/4.0/, as recorded in the arXiv licence metadata for this exact version and shown on the abstract page https://arxiv.org/abs/2605.27001v1. Full licence notice: \texttt{licenses/arxiv-2605.27001-CC-BY-4.0.txt}.\par\smallskip

\noindent\textit{Editorial scope.} Counted unit: the article's theorem thm2.7 (source0.tex lines 188--191). The article's complete original proof of it is delivered with it (source0.tex lines 192--202, one \texttt{proof} environment of 11 lines). No mathematics is written, repaired or re-derived: the statement and the proof are the article's own lines, in the article's own notation, and no retained line is substituted or re-flowed.  No further source-local context is needed: every cross-reference of every retained line resolves inside the two delivered spans.  Nothing was omitted from any retained span, so the package carries no omission record.  No retained line contains a citation, so the wrapper carries no bibliography and no bibliography entry.  No retained line contains a figure environment, an image-inclusion command or drawing code, so no image is delivered and the package declares no asset dependency.  Editorial formatting only, in addition: an AMS \texttt{amsart} preamble that restates the article's own declarations and macros used by the retained lines, namely the environments \texttt{theorem} on separate counters, together with the standard packages those lines need.  The retained environments print in this document as Theorem 1 in order of appearance, whereas the article prints the same items with the numbering of its own full document.  The complete native source of the article as distributed in the arXiv e-print for this exact version is bundled unchanged as \texttt{sources/201461/source0.tex}.
\begin{theorem}\label{thm2.7} 
Let $\{Q_i\}_{i=0}^n$ be $n+1$ hypersurfaces of $\P^n(\C)$ in general position of the same degree $p$ and $Q$ a hypersurface given by $Q=\sum_{i=0}^na_iQ_i^d$, where $a_i\ (0\le i\le n)$ are nonzero constants and $d$ is a positive integer. Let $f$ be a holomorphic mapping of $\Delta_{s,\infty}$ into $\P^n(\C)$ with a reduced representation $F=(f_0,\ldots,f_n)$ such that $Q_0^d(F),\ldots,Q_n^d(F)$ are linearly independent. Then, we have
$$ \biggl \|\ \left(1-\frac{n(n+1)}{d}\right)T_f(r,s)\le \frac{1}{pd}N^{[n]}_{Q(f)}(r,s)+o(T_f(r,s))+O(\log r).$$
\end{theorem}

\begin{proof}
Let $\tilde f$ be the holomorphic curve with the reduced representation $(Q_0^d(F),\ldots,Q_n^d(F))$. Then $\tilde f$ is linearly nondegenerate. By the second main theorem for hyperplanes, we have
\begin{align*}
\bigl\|\ T_{\tilde f}(r,s)&\le\sum_{i=0}^nN^{[n]}_{Q_i^d(f)}(r,s)+N^{[n]}_{Q(f)}(r,s)+o(T_{\tilde f}(r,s))+O(\log r)\\
&\le n\sum_{i=0}^nN_{Q_i(f)}(r,s)+N^{[n]}_{Q(f)}(r,s)+o(T_{\tilde f}(r,s))+O(\log r)\\
&\le n(n+1)pT_f(r,s)+N^{[n]}_{Q(f)}(r,s)+o(T_{\tilde f}(r,s))+O(\log r).
\end{align*}
Since $T_{\tilde f}(r)=dpT_f(r,s)+O(1)$, the above inequality implies that
$$ \biggl \|\ \left(1-\frac{n(n+1)}{d}\right)T_f(r,s)\le \frac{1}{pd}N^{[n]}_{Q(f)}(r,s)+o(T_f(r,s))+O(\log r).$$
The theorem is proved.
\end{proof}

\end{document}
